% !TeX root = ../navier-stokes-manuscript.tex
\subsection{Velocity, acceleration, jerk}\label{sub:BodyFiniteTimeResponseEscalation}

Acceleration is already in the \NavierStokes{} equation:
\begin{equation}\label{eq:MaterialAccelerationEquationBody}
  D_tu
  :=\partial_tu+(u\cdot\nabla)u
  =-\nabla p+\nu\Delta u.
\end{equation}
Jerk differentiates that response once more.
Because the material derivative itself contains transport, the next response already mixes time differentiation with spatial variation, pressure, and viscosity.

To expose the temporal rows at a fixed position, set
\begin{equation}\label{eq:TemporalTowerDefinitionBody}
  V_n:=\partial_t^nu,
  \qquad
  Q_n:=\partial_t^np,
  \qquad
  n\in\Nzero.
\end{equation}
Thus \(D_tu=V_1+(V_0\cdot\nabla)V_0\); \(V_1\) is not the material acceleration.
Solving for \(V_1\) and differentiating at fixed \(x\) gives
\[
\begin{aligned}
  V_1={}&
  \nu\Delta V_0
  -(V_0\cdot\nabla)V_0
  -\nabla Q_0,
  \\
  V_2={}&
  \nu\Delta V_1
  -(V_1\cdot\nabla)V_0
  -(V_0\cdot\nabla)V_1
  -\nabla Q_1,
  \\
  V_3={}&
  \nu\Delta V_2
  -(V_2\cdot\nabla)V_0
  -2(V_1\cdot\nabla)V_1
  -(V_0\cdot\nabla)V_2
  -\nabla Q_2,
  \\
  V_4={}&
  \nu\Delta V_3
  -(V_3\cdot\nabla)V_0
  -3(V_2\cdot\nabla)V_1
  -3(V_1\cdot\nabla)V_2
  -(V_0\cdot\nabla)V_3
  -\nabla Q_3.
\end{aligned}
\]
The coefficients are the binomial coefficients generated when the time derivatives split between the transporting and transported velocities.
At order \(n\),
\begin{equation}\label{eq:BodyTemporalVelocityRecurrence}
  V_{n+1}
  +\sum_{a=0}^{n}\binom{n}{a}(V_a\cdot\nabla)V_{n-a}
  +\nabla Q_n
  =\nu\Delta V_n,
  \qquad
  \nabla\cdot V_n=0.
\end{equation}

Pressure advances through the same rows.
With the decaying whole-space normalization,
\begin{equation}\label{eq:PressureNormalization}
  -\Delta p=\partial_i\partial_j(u_i u_j),
  \qquad
  p=R_iR_j(u_i u_j),
\end{equation}
and hence
\begin{equation}\label{eq:BodyTemporalPressureRecurrence}
  Q_n
  =R_iR_j
  \sum_{a=0}^{n}\binom{n}{a}
  (V_a)_i(V_{n-a})_j.
\end{equation}
The smooth datum and these two recurrences generate every finite initial temporal row.
Repeated integration can now make the Zeno demand exact: if one fixed \(V_n\) remained uniformly bounded in \(L^3\) on the finite interval, then the lower rows and finally \(u\) would remain uniformly bounded there as well.
Thus
\begin{equation}\label{eq:TemporalTowerMustEscape}
  T_*<\infty
  \quad\Longrightarrow\quad
  \sup_{0<t<T_*}\norm{V_n(t)}_3=\infty
  \qquad
  \text{for every fixed }n\in\Nzero.
\end{equation}
Different rows may lose control at different rates and in different regions.

The recurrence does not remain a temporal calculation.
Every row contains \(\Delta V_n\), spatial differentiation inside transport, and a pressure row reconstructed from the spatial field.
For \(n\in\Nzero\) and \(\alpha\in\Nzero^3\), set
\begin{equation}\label{eq:BodyMixedJetDefinition}
  J_{n,\alpha}:=\partial_t^n\partial_x^\alpha u,
  \qquad
  \Pi_{n,\alpha}:=\partial_t^n\partial_x^\alpha p.
\end{equation}
Applying \(\partial_x^\alpha\) to \eqref{eq:BodyTemporalVelocityRecurrence} gives
\begin{align}
  J_{n+1,\alpha}
  &+
  \sum_{a=0}^{n}
  \sum_{\beta\leq\alpha}
  \binom{n}{a}\binom{\alpha}{\beta}
  (J_{a,\beta}\cdot\nabla)
  J_{n-a,\alpha-\beta}
  +\nabla\Pi_{n,\alpha}
  \notag\\
  &=
  \nu\Delta J_{n,\alpha},
  \qquad
  \nabla\cdot J_{n,\alpha}=0.
  \label{eq:BodyMixedJetRecurrence}
\end{align}
The pressure coordinate satisfies
\begin{equation}\label{eq:BodyMixedPressureRecurrence}
  -\Delta\Pi_{n,\alpha}
  =
  \partial_x^\alpha\partial_i\partial_j
  \sum_{a=0}^{n}\binom{n}{a}
  (V_a)_i(V_{n-a})_j.
\end{equation}
Since \(u_0\in\mathscr S_\sigma(\R^3)\), finite-order Sobolev multiplication and boundedness of the Riesz transforms close the induction at \(t=0\):
\begin{equation}\label{eq:BodyInitialPressureCompleteMixedJet}
  J_{n,\alpha}(0)\in H_\sigma^m(\R^3),
  \qquad
  \Pi_{n,\alpha}(0)\in H^m(\R^3)
\end{equation}
for every finite \(n\), \(\alpha\), and \(m\).
This initializes the Tower; it does not control it near \(T_*\).

The object forced by the singularity is now visible.
Moving downward increases temporal response order.
Moving across increases spatial derivative order.
The derivative Tower is a time--space grid.
